%قبل از شروع مسئله ترجیحاً چیزی ننویسید. اما اگر خواستید بنویسید، در قالب لیستی (فقرات و شمارش) بنویسید.
%برای اضافه کردن نمره، بعدی از شروع‌مسئله، [x] را اضافه کنید و به جای x  نمره مورد نظر را قرار دهید. برای زیرمسئله‌ها هم به همین صورت است.
%\شروع‌مسئله[10]
\شروع‌مسئله

% صورت اصلی مسئله را در این‌جا بنویسید.

یگ گراف جهت‌دار در اختیار دارید که یکی از رئوس آن به عنوان شروع و راس دیگری به عنوان پایان در نظر گرفته شده است.


آ) توضیح دهید چگونه با استفاده از الگوریتم‌های جست‌وجوی عمق اول و جست‌وجوی سطح اول می‌توان گفت از راس شروع به راس پایان مسیر وجود دارد یا خیر. برای هر دو الگوریتم توضیح دهید.


ب) یک بار با استفاده از جست‌وجوی عمق اول و یک بار هم با استفاده از جست‌وجوی سطح اول، مسیر از راس شروع به راس پایان را در صورت وجود پیدا کنید.


پ) چگونه می‌توان تمامی مسیرهای موجود از شروع به پایان را پیدا کرد؟

\solution{
آ) برای فهمیدن این که آیا مسیری وجود دارد یا خیر، این دو الگوریتم تفاوت چندانی ندارند. کافی است آن‌ها را از راس شروع اجرا کنیم و هر راسی که پیمایش می‌شود را علامت‌گذاری کنیم. در نهایت باید چک کنیم راس پایان علامت‌گذاری شده است یا خیر.


ب)  با استفاده از جست‌وجوی سطح اول، کافی است الگوریتم را از راس شروع اجرا کنیم و برای هر راسی، نگه داریم که از کدام راس واردش شدیم. وقتی به راس پایان رسیدیم، پدران‌اش را صدا می‌زنیم تا به راس شروع برسیم و مسیر مشخص شده است. راه حل با استفاده از جست‌وجوی عمق اول به این صورت است که یک آرایه از گره‌ها در نظر می‌گیریم که در ابتدا فقط راس شروع در آن قرار داده شده است. الگوریتم را از راس شروع اجرا می‌کنیم، وقتی یک راس گسترده می‌شود، آن را
visited
می‌کنیم و به آرایه اضافه می‌کنیم. زمانی که راس
finish
می‌شود، آن را 
unvisited 
می‌کنیم و از آرایه حذف می‌کنیم. وقتی به راس پایان رسیدیم، گره‌های حاضر در لیست مسیری از راس شروع به راس پایان هستند و فرایند متوقف می‌شود.


پ) برای این کار از الگوریتم جست‌وجوی عمق اول استفاده می‌کنیم. همان الگوریتمی که در قسمت قبل گفته شد برای این قسمت هم کاراست با این تفاوت که در انتها، فرایند متوقف نمی‌شود و تا جایی ادامه می‌یابد که همه‌ی راس‌های گراف پیمایش شوند. هر بار در این روند به راس پایان رسیدیم، اعضای حاضر در آرایه چاپ شده و آرایه خالی می‌شود.

}



%%%%%%%%%%%%%%%%%%%%%%%%%%%%%%%%%%%%%%%%%%%%%%%%%%%%%
% برخی دستورات مفید برای تمرین‌ها. %
%%%%%%%%%%%%%%%%%%%%%%%%%%%%%%%%%%%%%%%%%%%%%%%%%%%%%

% فرمول‌های نمونه:
%$f(x) = \Theta(x)$
%$T(x, y) \leq \log(x) + 2^y + \binom{2n}{n}$
% {\tt code\_function}


% لیست موارد بدون شماره‌گذاری:
%\شروع{فقرات}
%    \فقره مورد اول 
%        \شروع{فقرات}
%            \فقره مورد اول داخلی 
%                \شروع{فقرات}
%                    \فقره اولین مورد خیلی داخلی 
%                    \فقره دومین مورد خیلی داخلی 
%                \پایان{فقرات}
%            \فقره مورد دوم داخلی 
%        \پایان{فقرات}
%    \فقره مورد دوم 
%\پایان{فقرات}

% لیست موارد شماره‌گذاری شده:
%\شروع{شمارش}
%    \فقره یک مورد 
%    \فقره یک مورد دیگر 
%    \فقره باز هم یک مورد دیگر
%\پایان{شمارش}

% برای نوشتن یک فرمول چند خطی که یک شماره بخوره. برای حالت‌های دیگر به سایت زیر مراجعه کنید:
% https://www.overleaf.com/learn/latex/Aligning_equations_with_amsmath
%\begin{align} 
%    \begin{split}
%        (x+y)^3 &= (x+y)^2(x+y) \\
%                &= (x^2+2xy+y^2)(x+y) \\
%                &= (x^3+2x^2y+xy^2) + (x^2y+2xy^2+y^3) \\
%                &= x^3+3x^2y+3xy^2+y^3
%    \end{split}                                 
%\end{align}

% نوشتن ماتریس و آرایه دوبعدی
%\begin{align}
%    A = 
%    \begin{bmatrix}
%        A_{11} & A_{21} \\
%        A_{21} & A_{22}
%    \end{bmatrix}
%\end{align}

% اضافه کردن تصویر با شرح. برای اضافه کردن تصویر بدون شرح، فقط دستور درج‌تصویر را استفاده کنید.
%\شروع{شکل}[h!]
%\تنظیم‌ازوسط
%  \درج‌تصویر[پهنا=0.5\پهنای‌متن]{img.jpg}
%\شرح{شرح تصویر}
%\برچسب{تصویر:۱}
%\پایان{شکل}
%\پایان{زیرمسائل}

%کد زیر را ببینید.
%%\lstinputlisting{source_filename.py}
%\begin{latin}
%\begin{lstlisting}
%# Import the modules
%import sys
%import random
%
%ans = True
%
%while ans:
%question = raw_input("Ask the magic 8 ball a question: (press enter to quit) ")
%
%answers = random.randint(1,8)
%
%if question == "":
%sys.exit()
%
%elif answers == 1:
%print "It is certain"
%
%elif answers == 2:
%print "Outlook good"
%
%elif answers == 3:
%print "You may rely on it"
%
%elif answers == 4:
%print "Ask again later"
%
%elif answers == 5:
%print "Concentrate and ask again"
%
%elif answers == 6:
%print "Reply hazy, try again"
%
%elif answers == 7:
%print "My reply is no"
%
%elif answers == 8:
%print "My sources say no"
%\end{lstlisting}
%\end{latin}
